% Korean edition of sections/s1_intro.tex.
\seclabel
\claim{\secnum{1}{sec:intro}\quad 공개 API가 없어 웹 클라이언트의 요청으로 서버를 측정했다}

\lead{대회는 웹 페이지로만 주문을 받으므로, 프로그램으로 매매하려면 브라우저를 대신할 래퍼가 필요하다.} Road to Fund Manager는 타임폴리오자산운용이 운영하는 모의투자 대회로, 참가자는 웹 페이지에서 종목코드와 목표 비중을 입력한다. 알고리즘 전략은 스스로 로그인하고, 잔고를 읽고, 바뀌는 포지션을 따라가고, 주문을 넣어야 한다. 래퍼가 이 통신을 맡으므로 전략 코드는 토큰, 소켓, 날짜 형식을 다루지 않는다. 래퍼는 페이지 자체가 보내는 HTTP 요청을 그대로 보내며, \ref{sec:design}장에서 브라우저로 페이지를 조작하는 방식과 비교한다.

\lead{주최측이 API를 공개하지 않으므로, 래퍼가 두는 모든 가정은 관찰에서 나와야 한다.} 인증 방식, 토큰 유효기간, 날짜 해석, 실시간 메시지의 내용이 모두 여기에 해당한다. 서버는 좀처럼 오류를 내지 않으므로 추측은 위험하다. \ref{sec:queries}절에서 보듯 잘못된 날짜를 보내면 HTTP 200과 함께 다른 날의 데이터가 돌아오므로, 잘못된 가정은 오류 없이 틀린 숫자를 만든다.

\lead{주최측은 두 가지 제한 안에서 자동 주문을 허용한다.} 2026년 11월 24일 우리 문의에 대한 답장[5]에서 주최측은 공식 API를 제공하지 않으며, 스크립트나 브라우저 자동화 도구(웹 드라이버)로 주문을 자동화하는 것을 금지하지 않는다고 밝혔다. 하루 100건 이상 주문하는 IP 주소와, 시스템에 큰 부하를 주는 순간적인 트래픽 폭주는 차단될 수 있다. 주최측은 요청 사이에 지연을 두라고 권했지만 수치는 주지 않았다. 답장 사본은 보관하고 있다. 우리는 부하 테스트나 요청 제한 테스트를 하지 않았다. 가장 많은 트래픽은 \ref{sec:transport}절의 시간 측정에서 나왔다. 순차 요청 5개(1회에 약 170\ms)를 1\,s 간격으로 10회 반복했다. 부하 기준은 알려지지 않았다.

\lead{래퍼가 쓰는 모든 요청 경로와 피드 메서드는 페이지의 JavaScript와 기록된 트래픽에서 얻었다.} 로그인한 뒤 사이트가 내려주는 JavaScript에서 API 경로와 실시간 피드 메서드 이름을 뽑고, 브라우저 개발자 도구로 WebSocket 메시지를 기록했다. 그다음 각 요청을 curl로 재현해 헤더, 쿠키, 응답 본문을 읽었다. curl로 다룰 수 없는 WebSocket은 클라이언트 라이브러리 없이 Python으로 SignalR 프로토콜을 직접 구현했다.

\lead{설명이 둘 이상 가능한 동작은 그 설명들을 구분하도록 설계한 요청으로 판정했다.} 경쟁 가설을 적어 두고, 가설마다 응답이 달라지는 요청을 골랐다. 델타가 오는 조건은 연결 2개에서 보낸 요청 11개로 후보 규칙 3개 중 1개만 남겼다(\ref{sec:feed}절). 서버의 주식 수 반올림은 1.76주에 해당하는 주문 1건으로 내림과 반올림을 구분했다(\ref{sec:oneshare}절).

\lead{보고서는 측정, 설계, 실서버 테스트 순서로 진행한다.} \ref{sec:layers}장은 프로토콜 5개 계층을 측정한다. 전송(\ref{sec:transport}), 인증(\ref{sec:auth}), 조회(\ref{sec:queries}), 실시간 피드(\ref{sec:feed}), 주문(\ref{sec:orders})이다. 앞의 4개는 주로 09:00 개장 전에 측정했고, 장중에 결과가 달라질 수 있는 항목은 개장 후 다시 확인했다. 주문은 장중에 측정했다. \ref{sec:design}장은 측정값을 래퍼의 목표, 구조, 운영 규칙으로 옮긴다. \ref{sec:live}장은 완성된 래퍼를 실서버에서 테스트한다. 시작, 피드, 주문(\ref{sec:wrappertests}), 주문이 0주로 내림되지 않게 막는 검사(\ref{sec:oneshare}), 대회 투자 가이드라인(\ref{sec:guidelines}), 그리고 이 테스트의 한계와 남은 질문(\ref{sec:validity})을 다룬다. \figref{setup}은 테스트 환경이다. 부록~\ref{app:glossary}는 용어를 정의하고, 부록~\ref{app:decisions}는 설계 결정과 근거를 나열한다.

% Korean edition of figs/setup.tex.
\fig{setup}{모든 테스트는 2026년 11월 24일 계정 하나와 클라이언트 하나로 했다}[이 보고서의 모든 측정에 쓴 패키지, 테스트 시간, 클라이언트, 계정, 서버, 부하]{%
{\tabfont\begin{tabularx}{\linewidth}{>{\RaggedRight\arraybackslash}p{16mm} >{\RaggedRight\arraybackslash}X}
\thd{항목} & \thd{값}\\
패키지 & \code{pytimefolio} v0.1.0, MIT 라이선스 [8]; \href{https://github.com/chlanstjr0301/pytimefolio}{\code{github.com/chlanstjr0301/pytimefolio}}\\\rowsep
시간 & 2026년 11월 24일 05:38\textasciitilde 10:34 KST. 시장은 09:00에 열렸다. 이 시간 동안 서버의 영업일(\code{currBizDay})이 11월 23일에서 24일로 바뀌었다\\\rowsep
클라이언트 & Windows 11의 WSL(Ubuntu). HTTP/2를 지원하는 curl 8. Python 3.14, h2를 포함한 httpx 0.28.1, websockets 17.4. PC 시계가 서버 \code{Date} 헤더보다 3\,s 빠름\\\rowsep
계정 & 참가자 계정 하나. 운용 중인 포트폴리오 하나(RFM 13회차, 2026년 10월 12일\textasciitilde 11월 30일)\\\rowsep
서버 & \code{https://contest.timefolio.net} (221.168.38.40). Let's Encrypt 인증서의 TLS 1.3. REST는 \code{/api} 아래, 실시간 피드는 \code{/hubs/main}\\\rowsep
부하 & 시간 측정: 약 170\ms{}에 순차 요청 5개, 1\,s 간격으로 10회. 부하·요청 제한 테스트는 없음. 실시간 피드 요청은 3\,s마다. 하루 주문: 새 주문 9건과 정지 7건(주최측 기준 100건)\\
\end{tabularx}}}{네트워크 시간은 이 클라이언트와 이 시간대에만 해당한다.}{DART Quant Team}

